% Author : Pierre-Henri Horrein
% Copyright : Telecom Bretagne, 2013

% Définition du document
\documentclass[t,9pt,pdftex]{beamer}
\usetheme{tb}
\usepackage{ucs}
\usepackage[utf8x]{inputenc}
\usepackage{textcomp}
\usepackage{graphicx}
\usepackage[french]{babel}
%\setbeamertemplate{navigation symbols}[horizontal]
\usepackage{multicol}

\graphicspath{{figs/}}


% Entre [], le titre raccourci pour le bas de page, 
% Entre {}, le titre complet pour la page de garde
\title[Modèle Télécom Bretagne]{Utilisation du modèle Latex de présentation pour
Télécom Bretagne}

% Nom de l'auteur
\author{Pierre-Henri Horrein}

% Date de la présentation (bas de page)
\date{14/03/2013}


\begin{document}

%---------- Titre -----------

% Création de la page de garde à partir des informations précédentes
\inserttitlepage


%---- Table des matières ----

% Une slide d'agenda
\begin{frame}[c]
	\frametitle{Agenda}
	% Plusieurs choix possibles (décommenter le choix sélectionné, commenter les
	% autres):
	% Uniquement les sections principales : 
	\tableofcontents[hideallsubsections]
	% Sections et sous sections : 
	%\tableofcontents
\end{frame}


\section{ Utilisation basique}

% On peut rappeler le plan à chaque début de section ou même de sous-section. La position de cette
% commande est importante : si on la met avant la première section, elle
% apparaîtra dans la présentation avant la première section

%\AtBeginSubsection[] {
%\AtBeginSection[] {
% \begin{frame}[c]
% \frametitle{Sommaire}
%   \small \tableofcontents[hideallsubsections,sectionstyle=show/shaded,subsectionstyle=hide/hide/hide ]
% \end{frame}
%}

\begin{frame}[c]
	\frametitle{Slide classique}
	\begin{itemize}
		\item on peut utiliser les niveaux classiques d'indentation
			\begin{itemize}
				\item niveau 2
					\begin{itemize}
						\item niveau 3
					\end{itemize}
			\end{itemize}
		\item les couleurs correspondent au modèle ppt
	\end{itemize}
\end{frame}

\section{Spécifités Beamer}

\begin{frame}[c]
	\frametitle{ Utilisation des blocs Latex}
	\begin{block}{Exemple de bloc normal}
		\begin{itemize}
			\item On peut utiliser un bloc normal, à fond vert
			\item ça met bien en valeur l'information essentielle de la slide
		\end{itemize}
	\end{block}
	\begin{exampleblock}{Exemple de bloc exemple}
		\begin{itemize}
			\item On peut utiliser un bloc exemple, à fond gris
			\item ça permet de séparer du rest
		\end{itemize}
	\end{exampleblock}
	\begin{alertblock}{Exemple de bloc alerte}
		\begin{itemize}
			\item On peut utiliser un bloc alerte, à fond marron
			\item ça permet d'avoir trois types de blocs
		\end{itemize}
	\end{alertblock}
\end{frame}

\section{Mises en page}

\begin{frame}[c]
	\frametitle{Deux colonnes}
	\begin{columns}
		\begin{column}[c]{.4\textwidth}
			% L'environnement figure est nécessaire ici
			\begin{figure}
				\centering
				\includegraphics[width=.9\textwidth]{Modele/bretagne_quadri}
			\end{figure}
		\end{column}
		\begin{column}[c]{.6\textwidth}
			\begin{itemize}
					\vfill
				\item on peut par exemple décrire l'image à gauche
					\vfill
				\item ou comparer deux colonnes
					\vfill
			\end{itemize}
		\end{column}
	\end{columns}
\end{frame}

\begin{frame}[c,containsverbatim]
\frametitle{ Code ou verbatim}
			\begin{verbatim}
				\centering
				\includegraphics[width=.9\textwidth]{Modele/bretagne_quadri}
			\end{verbatim}
			\begin{itemize}
					\vfill
				\item on peut aussi écrire du code ou du \verb+verbatim+
					\vfill
				\item au lieu des images
					\vfill
			\end{itemize}
		\end{frame}




\end{document}
